% Draft status: V0
% Evidence status: checked where possible; TODOs mark missing evidence.
\section{Datasets}

\noindent\textbf{GeoDE.}
GeoDE is a geographically diverse object-recognition dataset containing 61,940 images from 40 classes and six world regions \citep{ramaswamy2023geode}. The paper uses a processed train/validation/test split prepared for this experiment. The local class mapping used by the result files contains 40 classes, matching the official class count. The region labels appearing in the processed results are Africa, Americas, EastAsia, Europe, SouthEastAsia, and WestAsia. GeoDE is useful in this paper because it directly evaluates whether object recognition remains stable across broad geographic regions.

\noindent\textbf{Dollar Street.}
Dollar Street is a supervised household-object dataset with 38,479 images, 289 household-item categories, 404 homes, 63 countries, and four regions \citep{rojas2022dollarstreet}. The official dataset includes metadata such as region, country, and monthly income. This paper does not use the full 289-category official label space. Instead, it uses a processed ImageNet-compatible experimental subset with 64 classes, as recorded in the local class mapping. The processed result files use four region groups: Africa, Americas, Asia, and Europe. Dollar Street is useful here because everyday household objects can vary strongly in visual form across geography and socioeconomic setting.

\input{revision_dataset_summary_table}

The two datasets play complementary roles. GeoDE is designed specifically around geographically diverse object recognition, while Dollar Street emphasizes household objects embedded in geographic and socioeconomic context. The distinction matters for interpretation: a model that looks regionally stable on one dataset may still show larger geographic disparity on the other. TODO: add local split sizes, per-region sample counts, and class-by-region counts before the final version.
